\Article{15}
In the preceding memoir, Art. 37 and following, we taught a method of determining the precision of observations a posteriori as closely as possible. Namely, if approximate values of \(\pi\) quantities have become known through observations of equal precision, and are compared with the values obtained by calculation from the most plausible values of the \(\rho\) elements on which those quantities depend, the squares of the differences must be added and their sum divided by \(\pi-\rho\). The resulting quotient may then be regarded as an approximate value of the square of the mean error belonging to this kind of observation. Whenever the observations have unequal precision, these rules need to be changed only in that the squares must be multiplied by the weights of the observations before they are added, and the mean error obtained in this way must refer to observations whose weight has been taken as unity.

Now, in the present treatment, that sum clearly agrees with the sum \(S\), and the difference \(\pi-\rho\) with the number of conditional equations \(\sigma\). Thus, for the mean error of observations whose weight \(=1\), we shall have the expression \(\BareRoot{\frac S\sigma}\); this determination will deserve greater confidence the larger the number \(\sigma\) is.

But it will be worthwhile to establish this independently of the preceding investigation as well. For this purpose it will be useful to introduce some new notation. Namely, let the following values of the indeterminates \(\xi,\eta,\zeta\) etc.,
\SourceLeaf{079}
\[
\xi=a,\quad\eta=b,\quad\zeta=c\quad\text{etc.}
\]
correspond to these values of \(x,y,z\) etc.,
\[
x=\alpha,\quad y=\beta,\quad z=\gamma\quad\text{etc.}
\]
so that we have
\[
\begin{aligned}
\alpha&=a[\alpha\alpha]+b[\alpha\beta]+c[\alpha\gamma]+\text{etc.}\\
\beta&=a[\alpha\beta]+b[\beta\beta]+c[\beta\gamma]+\text{etc.}\\
\gamma&=a[\alpha\gamma]+b[\beta\gamma]+c[\gamma\gamma]+\text{etc.}
\end{aligned}
\]
etc. Likewise, we shall suppose that to the values
\[
\xi=a',\quad\eta=b',\quad\zeta=c'\quad\text{etc.}
\]
there correspond these,
\[
x=\alpha',\quad y=\beta',\quad z=\gamma'\quad\text{etc.}
\]
and to these,
\[
\xi=a'',\quad\eta=b'',\quad\zeta=c''\quad\text{etc.}
\]
the following,
\[
x=\alpha'',\quad y=\beta'',\quad z=\gamma''\quad\text{etc.}
\]
and so on.

With these assumptions, combining equations (4) and (9) gives
\[
\begin{aligned}
A&=\alpha e+\alpha'e'+\alpha''e''+\text{etc.}\\
B&=\beta e+\beta'e'+\beta''e''+\text{etc.}\\
C&=\gamma e+\gamma'e'+\gamma''e''+\text{etc.}
\end{aligned}
\]
etc. Thus, since we have \(S=\mathfrak A A+\mathfrak B B+\mathfrak C C+\text{etc.}\), it is clear that
\[
\begin{aligned}
S={}&(ae+a'e'+a''e''+\text{etc.})(\alpha e+\alpha'e'+\alpha''e''+\text{etc.})\\
&+(be+b'e'+b''e''+\text{etc.})(\beta e+\beta'e'+\beta''e''+\text{etc.})\\
&+(ce+c'e'+c''e''+\text{etc.})(\gamma e+\gamma'e'+\gamma''e''+\text{etc.})+\text{etc.}
\end{aligned}
\]
\Article{16}
We may regard taking the observations, through which we obtain values of the quantities \(v,v',v''\) etc. affected by random errors \(e,e',e''\) etc., as
\SourceLeaf{080}
an experiment which cannot tell us the magnitude of each error committed, but which, by applying the rules explained above, supplies the value of the quantity \(S\), which by the formula just found is a given function of those errors. In such an experiment, some of the random errors may of course turn out larger and others smaller; but the more errors there are, the greater the hope that the value of the quantity \(S\) in an individual experiment will deviate little from its mean value. The essential point is therefore to establish the mean value of the quantity \(S\). By the principles presented in the preceding memoir, which it would be superfluous to repeat here, we find this mean value
\[
\begin{aligned}
&=(a\alpha+b\beta+c\gamma+\text{etc.})mm\\&\qquad+(a'\alpha'+b'\beta'+c'\gamma'+\text{etc.})m'm'\\
&\qquad+(a''\alpha''+b''\beta''+c''\gamma''+\text{etc.})m''m''+\text{etc.}
\end{aligned}
\]
\Needspace{20\baselineskip}

Denoting the mean error of observations whose weight \(=1\) by \(\mu\), so that \(\mu\mu=pmm=p'm'm'=p''m''m''\) etc., the expression just found can be written as follows:
\[
\begin{aligned}
&\left(\frac{a\alpha}{p}+\frac{a'\alpha'}{p'}+\frac{a''\alpha''}{p''}\ \text{etc.}\right)\mu\mu
\\&\qquad+\left(\frac{b\beta}{p}+\frac{b'\beta'}{p'}+\frac{b''\beta''}{p''}+\text{etc.}\right)\mu\mu\\
&\qquad+\left(\frac{c\gamma}{p}+\frac{c'\gamma'}{p'}+\frac{c''\gamma''}{p''}+\text{etc.}\right)\mu\mu+\text{etc.}
\end{aligned}
\]
But the sum \(\frac{a\alpha}{p}+\frac{a'\alpha'}{p'}+\frac{a''\alpha''}{p''}+\text{etc.}\) is found to be
\[
=[aa].[\alpha\alpha]+[ab].[\alpha\beta]+[ac].[\alpha\gamma]+\text{etc.}
\]
and therefore \(=1\), as is easily understood from the connection between equations (6) and (7). Likewise,
\[
\begin{aligned}
\frac{b\beta}{p}+\frac{b'\beta'}{p'}+\frac{b''\beta''}{p''}+\text{etc.}&=1\\[5pt]
\frac{c\gamma}{p}+\frac{c'\gamma'}{p'}+\frac{c''\gamma''}{p''}+\text{etc.}&=1
\end{aligned}
\]
and so on.

Hence, finally, the mean value of \(S\) is \(=\sigma\mu\mu\), and thus, insofar as it is permissible to adopt the random value of \(S\) as its mean value, we shall have \(\mu=\BareRoot{\frac S\sigma}\).

\clearpage
